% Supplementary material -- own numbering (Figures S1-S5, Tables S1-S2). Build: latexmk -pdf supplementary.tex
\documentclass[preprint,12pt]{elsarticle}
\usepackage{amsmath}
\usepackage{graphicx}
\usepackage{booktabs}
\usepackage{siunitx}
\usepackage{xcolor}
\usepackage[hidelinks]{hyperref}
\graphicspath{{figures/}}
\newcommand{\todo}[1]{\textcolor{red}{[TODO: #1]}}
\newcommand{\rdtversion}{v2}
\renewcommand{\thefigure}{S\arabic{figure}}
\renewcommand{\thetable}{S\arabic{table}}
\renewcommand{\thesection}{S\arabic{section}}
% auto-generated by rdt.paper_tables.numbers() from v2 result files; do not edit
\newcommand{\nLHS}{2000}
\newcommand{\nSaltelli}{4096}
\newcommand{\nMC}{2048}
\newcommand{\speedup}{4658}
\newcommand{\surrRMSETwo}{0.11}
\newcommand{\surrRMSElogtr}{0.0057}
\newcommand{\PICPtwo}{0.957}
\newcommand{\calFgt}{0.323}
\newcommand{\calLq}{0.456}
\newcommand{\calFhtg}{1.83}
\newcommand{\calChi}{34.9}
\newcommand{\calRedChi}{1.46}
\newcommand{\alphaIplant}{1408}
\newcommand{\alphaIxf}{348}
\newcommand{\xfRMSET}{14.0}
\newcommand{\xfRMSEinc}{0.015}
\newcommand{\SobolFiringTwo}{0.35}
\newcommand{\SobolSCTwo}{0.28}
\newcommand{\SobolEAtwo}{0.19}
\newcommand{\SobolPlogtr}{0.16}
\newcommand{\SobolFiringSlip}{0.54}
\newcommand{\SxDaily}{0.991}
\newcommand{\SxRenew}{0.983}
\newcommand{\SxSClow}{1.16}
\newcommand{\SxSChigh}{0.83}
\newcommand{\SxAgeSlip}{2.25}
\newcommand{\SxAgeTout}{1.88}
\newcommand{\SxAgeSlipYone}{0.71}
\newcommand{\TwoAgeStart}{1137}
\newcommand{\TwoAgeEnd}{1164}
\newcommand{\SxAvgErr}{0.982}
\newcommand{\evMild}{2.4}
\newcommand{\evSevere}{6.9}
\newcommand{\evHot}{11.1}
\newcommand{\evHotShare}{2.7}
\newcommand{\nPareto}{200}
\newcommand{\nFeasible}{199}
\newcommand{\nDominating}{37}
\newcommand{\kneeHtwo}{+16}
\newcommand{\kneeLife}{0.28}
\newcommand{\kneeFuel}{-18}
\newcommand{\isoLife}{0.083}
\newcommand{\isoFuel}{-16}
\newcommand{\isoTwo}{1109}
\newcommand{\baseTwo}{1147}
\newcommand{\maxHtwoLife}{4.3}
\newcommand{\maxHtwoGain}{+30}
\newcommand{\steamMinLifeZero}{0.059}
\newcommand{\steamMinLifeZeroFuel}{+3.2}
\newcommand{\steamMinFuelZeroSC}{2.66}
\newcommand{\steamMinFuelOneSC}{4.00}
\newcommand{\uqTwoWidth}{42}
\newcommand{\uqLogtrLo}{7.04}
\newcommand{\uqLogtrHi}{8.73}
\newcommand{\uqKneeLo}{0.20}
\newcommand{\uqKneeHi}{0.64}
\newcommand{\uqAgeLo}{1.84}
\newcommand{\uqAgeHi}{8.97}
\newcommand{\shareCreep}{62}
\newcommand{\shareKinTwo}{51}
\newcommand{\shareHT}{5}
\newcommand{\robAge}{0.91}
\newcommand{\robKnee}{0.81}
\newcommand{\robSC}{0.84}
\newcommand{\robFlexNom}{0.63}
\newcommand{\robFlexVar}{0.001}
\newcommand{\lqvarLoadLow}{1.71}
\newcommand{\pipelineMin}{34}

\journal{International Journal of Hydrogen Energy}

\begin{document}
\begin{frontmatter}
\title{Supplementary material to: A physics-informed digital twin for tube-life-aware operation of industrial steam methane reformers under uncertainty}
\author[su]{Izbassar Orynbassar\corref{cor1}}
\author[sua]{\todo{Second author name}}
\cortext[cor1]{Corresponding author.}
\address[su]{School of Information Studies, Syracuse University, Syracuse, NY, USA}
\address[sua]{Safi Utebayev Atyrau University of Oil and Gas, Atyrau, Kazakhstan}
\end{frontmatter}

\section{Digitised verification data}

The verification targets of Section 3.1 are the axial profiles of Fig.~3 of Xu and Froment (1989, Part II). The figure was
digitised curve by curve \todo{name the tool (e.g. WebPlotDigitizer), the image resolution and the axis-calibration procedure} and
stored as one tidy CSV per curve in \texttt{data/literature\_validation/digitized/} with the columns \texttt{z\_m}, value and a
\texttt{source} annotation. Curves: methane conversion, CO$_2$ conversion, process-gas temperature, inner-wall temperature,
outer-wall temperature and total pressure, over the heated length of \SI{11.12}{m}; the remaining length to \SI{12}{m} is treated as
adiabatic. The outer-wall curve is interpolated by a shape-preserving cubic (PCHIP) and used as the boundary condition; the
digitising uncertainty is estimated at \todo{quote the estimated reading uncertainty in K and in conversion} from repeat readings of
the axis ticks. The feed and geometry are those of Table~2 of the source, with a bed voidage of 0.526 and an equivalent ring diameter
of \SI{9.44}{mm} fitted to the pressure curve \todo{confirm that fitting the particle diameter to the pressure drop is acceptable}.

\section{Additional figures}

\begin{figure}[htbp]
\centering
\includegraphics[width=0.6\textwidth]{fig00_equilibrium_conversion}
\caption{Equilibrium methane conversion at \SI{25}{bar} and steam-to-carbon 3 from Cantera (GRI-Mech 3.0), used as a sanity check of
the thermodynamic basis. Source: \texttt{notebooks/00\_sanity\_check.ipynb}.}
\end{figure}

\begin{figure}[htbp]
\centering
\includegraphics[width=\textwidth]{fig01_kinetics}
\caption{Xu--Froment kinetics as implemented: (a) rate coefficients $k_1$--$k_3$ against reciprocal temperature with the Table~4
values of the source for $k_1$, and (b) rates $r_1$--$r_3$ at \SI{25}{bar} for a feed of CH$_4$:H$_2$O:H$_2$ = 1:3:0.1. Source:
\texttt{data/kinetics/xu\_froment\_1989.yaml}.}
\end{figure}

\begin{figure}[htbp]
\centering
\includegraphics[width=\textwidth]{fig02_reactor_profiles}
\caption{Stand-alone tube model with a linear outer-wall profile from 1000 to \SI{1150}{K} on the Xu--Froment geometry: (a)
conversion and CO$_2$ yield, (b) wall and gas temperatures, (c) pressure.}
\end{figure}

\begin{figure}[htbp]
\centering
\includegraphics[width=\textwidth]{fig04_latham_baseline}
\caption{Coupled model with the literature (uncalibrated) parameters on all four Latham plant cases, including Plant C1 which is
excluded from calibration. Source: \texttt{latham2008\_plant\_cases.csv}, \texttt{latham\_fit.yaml} (baseline).}
\end{figure}

\begin{figure}[htbp]
\centering
\includegraphics[width=0.6\textwidth]{fig10_scenarios_timeseries}
\caption{One week (days 140--147) of the renewable-following scenario S3: load, firing factor solved by the outlet-temperature
controller, hot-spot temperature and hourly damage rate. Source: \texttt{data/scenarios/hourly\_v2\_S3\_renewable.csv.gz}.}
\end{figure}

\section{Operating space}

\begin{table}[htbp]
\centering
\caption{Operating inputs and their ranges (v2 operating space); base = calibrated Plant A. Source:
\texttt{data/design/parameter\_ranges\_v2.csv}.}
\label{tab:ranges}
\small
\begin{tabular}{llrrrp{5.5cm}}
\toprule
Input & Unit & Min & Max & Base & Justification (abridged) \\
\midrule
steam-to-carbon & -- & 2.5 & 4.0 & 2.89 & carbon-formation limit below $\approx$2.5; 4.0 upper for hydrogen plants \\
load (feed per tube) & -- & 0.60 & 1.10 & 1.00 & turndown 50--60\,\%; 110\,\% short-term overload \\
specific firing factor & -- & 0.85 & 1.15 & 1.00 & burner control range \\
inlet temperature & K & 800 & 900 & 884.6 & Plant A 611 \si{\celsius}; cracking in preheat coils above $\approx$920 K \\
inlet pressure & bar & 25 & 33 & 30.06 & industrial 20--40 bar \\
catalyst activity & -- & 0.10 & 0.30 & 0.20 & Latham fitted 0.20 for a mid-campaign catalyst \\
excess air & \% & 5 & 25 & 21.6 & upper bound raised from 20 to 25\,\% so that the calibrated Plant A value lies inside the space \\
\bottomrule
\end{tabular}
\end{table}

\section{Version note}

All results in the main text are from pipeline version \rdtversion, which differs from the exploratory version v1 in the excess-air
range (5--25\,\% instead of 5--20\,\%), in the use of the calibrated excess air as the base of the scenarios and the optimisation
(instead of 20\,\%) and in the Monte Carlo sample size (\nMC{} instead of 4000). A line-by-line comparison of the two versions is
provided in \texttt{data/design/v1\_vs\_v2\_diff.md}; the largest effect of the change is a base hot spot about 3 K cooler and hence
a longer base life, which shifts every relative life metric by the same denominator.

\end{document}
